%% =====================================================================
%%  附录 A　常见错误码对照表
%%  源文件：附录A-错误码对照表.md
%%  编号体系独立于正文（附录 A / A.1 …）
%% =====================================================================
\chapter{常见错误码对照表}

\applicable{你收到一句话，不知道它到底是什么意思。}

\section{使用说明}

值得注意的是，人类接口的错误提示有一个特征，本手册在 0.3 里已经提过一次：

\begin{mdquote}
\textbf{含糊度与问题严重程度正相关。}对方说得越轻描淡写，问题通常越大。
\end{mdquote}

下表是样本库里出现频率最高的一批返回，按收到的话排列。\textbf{它不是字典}——只描述“最可能”，不描述“一定”。使用时请先对照基线（第 5 章）。

\begin{manstable}{P{3.6cm} P{3.0cm} L L}
\tbh{收到的话} & \tbh{类状态码} & \tbh{实际含义} & \tbh{处理} \\
\midrule
\textbf{“嗯”} & 100 Continue & 收到了，继续发 & 继续。且 20—40 秒后再给一次 \\
\textbf{“在吗”} & 102 Processing & 我要发一个请求，先握个手 & 回“在”。但\hciSee{第1章}，这是坏请求的开头 \\
\textbf{“好的”} & 200 OK & 同意，无附加信息 & 正常推进 \\
\textbf{“好的。”} & 200 OK（带签名） & 同意，但态度被收回了 & 看基线偏离，可能是信号 \\
\textbf{“挺好的”} & 200 OK（空包体） & 分两种：真好 / 很不好但不想说 & 看语气和延迟的偏离 \\
\textbf{“随便”} & 200 OK（body 缺失） & 我在乎，但我不想先说，也不想负责 & 给两三个具体选项，不要真的随便 \\
\textbf{“都行”} & 200 OK & 有不行，但我不打算说 & 同上 \\
\textbf{“你看着办吧”} & 200 OK & 我有倾向，而且你办错我会记住 & 给出选项并请对方挑 \\
\textbf{“随便你”} & 200 OK（带警告） & 我不高兴，且我不打算讨论 & \textbf{不要真的随便} \\
\textbf{“哦”} & 204 No Content & 收到，并且不高兴 & 停一下。不要追问 \\
\textbf{“呵呵”} & 403 Forbidden & 拒绝继续这个话题 & 换话题。不要解释 \\
\textbf{“没事”} & 200 OK，包体为空 & \hciSee{第6章}，至少四种含义 & 不追问，给一条不要求回答的通道 \\
\textbf{“算了”} & 499 Client Closed Request & 我本来想说，是你刚才那句让我不说了 & 停一下，说“你说，我听着” \\
\textbf{“你要这么想，我也没办法”} & 500 Internal Server Error & 我这边出问题了，且我不打算修 & \hciSee{第7章}。这是重伤信号 \\
\textbf{“我们过段时间再说”} & 202 Accepted & 已受理，暂不处理 & \textbf{要一个回程时间}，\hciSee{第12章} \\
\textbf{长时间未回复} & 408 Request Timeout & 未响应 & 对照这个人的基线，不要读绝对值 \\
\textbf{“你想太多了”} & 400 Bad Request & 你的请求格式不对，我不处理 & 注意：这是\textbf{它}的错误码，不是你的 \\
\textbf{“我早就跟你说过”} & 409 Conflict & 我在追认一次旧账 & \hciSee{第9章}，这句话会作废前面全部 \\
\textbf{已读不回} & 102 Processing（无后续） & 收到了，暂时无法处理 & \hciSee{第10章}。不要在本地计算返回 \\
\bottomrule
\end{manstable}

\section{使用本表的三条限制}

\textbf{一、本表描述的是“最可能”，不是“一定”。}同一句话在不同的人、不同的关系、不同的时段，含义可以完全不同。

\textbf{二、本表在基线缺失时不可用。}你和一个人认识不足一个月，请不要查表。

\textbf{三、本表会过期。}人的行为会在版本间静默变更（见附录 B）。半年后请重新校准。

%% 附录末「页脚交互条」由 hci-manual.cls 自动挂接，无需手写。
